\documentclass[aspectratio=169]{beamer}

% ============================================================
% Preámbulo autocontenido (no depende de preambulo.tex)
% ============================================================
\usetheme{Boadilla}
\usecolortheme{default}

\usepackage[utf8]{inputenc}
\usepackage[spanish,mexico]{babel}
\usepackage[T1]{fontenc}
\usepackage{amsmath,amssymb}
\usepackage{listings}
\usepackage{xcolor}
\usepackage{hyperref}
\usepackage{tikz}

\setbeamertemplate{navigation symbols}{}
\setbeamertemplate{footline}[frame number]
\setbeamercolor{block title}{bg=structure.fg!20,fg=black}

% Estilo para código R
\lstdefinestyle{Rstyle}{
  language=R,
  basicstyle=\ttfamily\small,
  keywordstyle=\color{blue},
  commentstyle=\color{gray},
  stringstyle=\color{teal},
  breaklines=true,
  showstringspaces=false,
  frame=single,
  backgroundcolor=\color{black!3}
}
\lstset{style=Rstyle}

\title[Introducción a R y RStudio]{Introducción a R y RStudio}
\subtitle{EST-11102 · Inferencia Estadística}
\author{C. Vladimir Rodríguez-Caballero}
\institute[ITAM]{Departamento Académico de Estadística, ITAM}
\date{\today}

\begin{document}

% ------------------------------------------------------------
\begin{frame}
\titlepage
\end{frame}

% ------------------------------------------------------------
\begin{frame}{Objetivo de esta sesión}
\begin{itemize}
  \item<1-> Familiarizarnos con \textbf{R} y \textbf{RStudio}, las herramientas de cómputo del curso.
  \item<2-> \textbf{No} es un curso de programación: es una introducción práctica para que puedan trabajar sin fricción.
  \item<3-> Todo el código del curso (ejemplos, tareas) estará disponible en R.
\end{itemize}
\end{frame}

% ------------------------------------------------------------
\begin{frame}{¿Por qué R?}
\begin{itemize}
  \item<1-> \textbf{Software libre} y de código abierto, mantenido por la comunidad estadística.
  \item<2-> Es el \textit{lenguaje estándar} en estadística académica y muy usado en econometría aplicada.
  \item<3-> Enorme ecosistema de paquetes: series de tiempo, econometría, finanzas, macro.
  \item<4-> Reproducibilidad: el mismo script genera siempre el mismo resultado.
\end{itemize}

\begin{block}<5->{Contexto económico}
Bancos centrales, organismos como el Banco Mundial y el FMI, y buena parte de la industria financiera usan R para análisis de datos y modelación.
\end{block}
\end{frame}

% ------------------------------------------------------------
\begin{frame}{R vs. RStudio: ¿cuál es la diferencia?}
\begin{columns}[T]
\begin{column}{0.48\textwidth}
\begin{exampleblock}{R}
\begin{itemize}
  \item El \textbf{motor}: el lenguaje y el intérprete.
  \item Se descarga de \href{https://cran.r-project.org}{CRAN}.
  \item Se puede usar solo, pero es incómodo.
\end{itemize}
\end{exampleblock}
\end{column}
\begin{column}{0.48\textwidth}
\begin{exampleblock}{RStudio}
\begin{itemize}
  \item El \textbf{IDE}: entorno de desarrollo.
  \item Se descarga de \href{https://posit.co/download/rstudio-desktop/}{posit.co}.
  \item Requiere tener R ya instalado.
\end{itemize}
\end{exampleblock}
\end{column}
\end{columns}

\vspace{0.5cm}
\pause
\centering
\textbf{Analogía:} R es el motor del coche; RStudio es el tablero, el volante y los asientos.
\end{frame}

% ------------------------------------------------------------
\begin{frame}{Instalación (orden importa)}
\begin{enumerate}
  \item<1-> Instalar \textbf{R} desde \url{https://cran.r-project.org}
  \item<2-> Instalar \textbf{RStudio Desktop} (versión gratuita) desde \url{https://posit.co/download/rstudio-desktop/}
  \item<3-> Abrir RStudio (no R directamente) para todo el trabajo del curso
\end{enumerate}

\vspace{0.3cm}
\pause
\begin{alertblock}{Nota}
Si actualizan R en el futuro, RStudio detecta la nueva versión automáticamente; no es necesario reinstalar RStudio.
\end{alertblock}
\end{frame}

% ------------------------------------------------------------
\begin{frame}{La interfaz de RStudio}
\centering
\begin{tikzpicture}[scale=0.85, every node/.style={font=\small}]
  \draw (0,0) rectangle (10,6);
  \draw (5,3) -- (5,6);
  \draw (0,3) -- (10,3);
  \draw (5,0) -- (5,3);

  \node at (2.5,5.4) {\textbf{Script (.R)}};
  \node[align=center] at (2.5,4.3) {donde se escribe\\ y guarda el código};

  \node at (7.5,5.4) {\textbf{Environment / History}};
  \node[align=center] at (7.5,4.3) {objetos creados\\ en la sesión};

  \node at (2.5,2.4) {\textbf{Consola}};
  \node[align=center] at (2.5,1.3) {donde el código\\ se ejecuta};

  \node[align=center] at (7.5,2.35) {\textbf{Files / Plots /}\\ \textbf{Packages / Help}};
  \node[align=center] at (7.5,0.7) {gráficas, paquetes\\ y documentación};
\end{tikzpicture}
\end{frame}

% ------------------------------------------------------------
\begin{frame}[fragile]{Consola vs. Script}
\begin{columns}[T]
\begin{column}{0.5\textwidth}
\textbf{Consola}
\begin{itemize}
  \item Ejecuta código \textit{inmediatamente}.
  \item No se guarda automáticamente.
  \item Útil para probar cosas rápido.
\end{itemize}
\end{column}
\begin{column}{0.5\textwidth}
\textbf{Script (.R)}
\begin{itemize}
  \item Se guarda como archivo.
  \item Se ejecuta línea por línea con \texttt{Ctrl+Enter} (\texttt{Cmd+Enter} en Mac).
  \item Es lo que usaremos siempre en el curso.
\end{itemize}
\end{column}
\end{columns}

\vspace{0.4cm}
\pause
\begin{lstlisting}
# Esto es un comentario
x <- 5        # asignacion
y <- 3
x + y         # se ejecuta y muestra 8
\end{lstlisting}
\end{frame}

% ------------------------------------------------------------
\begin{frame}[fragile]{Conceptos básicos I: asignación y tipos}
\begin{lstlisting}
# Asignacion (preferir <- sobre =)
pib <- 1200000
tasa_interes <- 0.075
pais <- "Mexico"
es_miembro_ocde <- TRUE

class(pib)             # "numeric"
class(pais)             # "character"
class(es_miembro_ocde)  # "logical"
\end{lstlisting}
\pause
\begin{itemize}
  \item R distingue mayúsculas/minúsculas: \texttt{PIB} $\neq$ \texttt{pib}.
  \item Nombres de objetos: sin espacios, no empezar con número.
\end{itemize}
\end{frame}

% ------------------------------------------------------------
\begin{frame}[fragile]{Conceptos básicos II: vectores y data frames}
\begin{lstlisting}
# Vector: la estructura mas basica
inflacion <- c(4.2, 4.5, 3.9, 4.1)
mean(inflacion)
sd(inflacion)

# Data frame: tabla de datos (como Excel)
datos <- data.frame(
  pais = c("Mexico", "Chile", "Colombia"),
  inflacion_2024 = c(4.2, 4.0, 6.1)
)
head(datos)
\end{lstlisting}
\end{frame}

% ------------------------------------------------------------
\begin{frame}[fragile]{Paquetes: extendiendo R}
\begin{itemize}
  \item<1-> R base ya trae mucho, pero los \textbf{paquetes} añaden funciones especializadas.
  \item<2-> Se instalan \textit{una vez}, se cargan \textit{cada sesión}.
\end{itemize}

\vspace{0.3cm}
\pause
\begin{lstlisting}
# Instalar (solo una vez por computadora)
install.packages("tidyverse")

# Cargar (cada vez que se abre RStudio)
library(tidyverse)
\end{lstlisting}

\pause
\begin{block}{Paquetes que usaremos en el curso}
\texttt{tidyverse}, \texttt{tseries}, \texttt{nortest}, entre otros para pruebas de hipótesis y validación de normalidad.
\end{block}
\end{frame}

% ------------------------------------------------------------
\begin{frame}{Organización: RStudio Projects}
\begin{itemize}
  \item<1-> Un \textbf{proyecto} (\texttt{.Rproj}) fija la carpeta de trabajo automáticamente.
  \item<2-> Evita el uso de \texttt{setwd()} con rutas absolutas (que no funcionan en la computadora de alguien más).
  \item<3-> \textbf{Recomendación:} un proyecto por curso o por tarea, con subcarpetas \texttt{datos/}, \texttt{scripts/}, \texttt{figuras/}.
\end{itemize}

\pause
\begin{alertblock}{Buena práctica}
Todo script del curso debe poder ejecutarse de principio a fin sin errores en una sesión nueva de R (``reproducibilidad'').
\end{alertblock}
\end{frame}

% ------------------------------------------------------------
\begin{frame}{Ayuda y recursos}
\begin{itemize}
  \item \texttt{?funcion} o \texttt{help(funcion)} dentro de R.
  \item Paquete \texttt{swirl}: tutoriales interactivos dentro de la consola.
  \item Libro gratuito \textit{R for Data Science} (Wickham \& Grolemund).
  \item \href{https://cran.r-project.org/web/views/Econometrics.html}{CRAN Task View: Econometrics} — paquetes especializados.
  \item Los scripts del curso estarán en el sitio del curso (GitHub Pages).
\end{itemize}
\end{frame}

% ------------------------------------------------------------
\begin{frame}{Lo que viene}
\begin{itemize}
  \item De aquí en adelante, cada tema del curso vendrá acompañado de código en R.
  \item El objetivo \textbf{no} es programar de forma elegante, sino \textbf{entender la estadística} usando la computadora como herramienta.
  \item Practiquen: instalar R y RStudio antes de la próxima clase.
\end{itemize}

\vspace{0.5cm}
\centering
\Large \textbf{¡Bienvenidos a R!}
\end{frame}

\end{document}
